\chapter{Background}
\label{sec:background}

This chapter introduces what the later chapters build on. It begins with the biology and the
measurements behind the data, from tumour heterogeneity through bulk and single-cell RNA sequencing to
the representation of single-cell data and the cell-line screens that supply the labels. It then turns
to previous work on predicting drug response, in bulk and at single-cell resolution, and to the two
computational approaches this work uses, single-cell foundation models and multiple-instance learning.
It closes with the objective of this work.

\section{Tumour heterogeneity and drug resistance}
\label{sec:bg-heterogeneity}

A tumour is not a uniform population of cells. Its malignant cells carry diverse genetic and epigenetic
material, on which selection and Darwinian evolution act \cite{mcgranahan2017}, and they differ further in
phenotype, which integrates genetic, epigenetic and environmental inputs \cite{marusyk2020}. Non-genetic mechanisms are
argued to shape tumour-cell phenotypes more than mutations do, and tumour cells show high epigenetic and
phenotypic plasticity \cite{marusyk2020}. This intratumour heterogeneity (ITH) is a key challenge in
cancer medicine and a key factor in therapeutic failure and drug resistance \cite{mcgranahan2017}, and it
has long been a barrier to understanding cancer \cite{castro2021}.

Its clinical consequence is resistance. Most advanced cancers, including those that first respond
strongly, eventually acquire resistance to targeted therapies and relapse, and this acquired resistance
is the direct consequence of pre-existing heterogeneity and of ongoing diversification during therapy,
which lets some tumour cells survive treatment \cite{marusyk2020}. Survival need not be genetic: cells can
pass into a reversible, drug-tolerant state that survives therapy without itself sustaining regrowth, and
genuine resistance can then arise from such cells by new mutations \cite{marusyk2020}. Single-cell
profiling has found resistance-associated states before treatment, for example in BRAF-mutant melanoma
\cite{castro2021}. In a multiple myeloma trial of 28 patients, the predicted response of the most
resistant transcriptional clone, weighted by its abundance, separated responders from non-responders
better than an abundance-weighted average over clones (area under the receiver operating characteristic
curve, AUROC, 0.83 against 0.59), although in cell lines the average performed better \cite{perception}.

Bulk measurements cannot see this structure. Genomic and transcriptomic profiling has traditionally
measured tumours in bulk, revealing aggregate profiles and masking important aspects of ITH
\cite{castro2021}, and the drug screens that provide response labels measure bulk populations as well
(Section~\ref{sec:bg-screens}). The cancer cell lines used in those screens are themselves heterogeneous:
across the cell lines of the SCP542 atlas, 12 expression programs were recurrently heterogeneous within
multiple lines, and co-existing subpopulations sorted from two head and neck cancer lines differed in
their sensitivity to drugs \cite{scp542}. In the same atlas, genetic subclones, inferred from
copy-number alterations (gains or losses of chromosome segments), were detectable in 26\,\% of the
\NLinesAtlas{} lines, and 39\,\% of the discrete expression clusters but only 8\,\% of the continuous
expression programs were associated with them \cite{scp542}.

\section{Bulk and single-cell RNA sequencing}
\label{sec:bg-rnaseq}

RNA sequencing (RNA-seq) profiles a transcriptome. In the standard workflow, RNA is extracted, enriched
for messenger RNA or depleted of ribosomal RNA, converted to complementary DNA (cDNA) and prepared as an
adaptor-ligated library, which is sequenced to 10--30 million short reads per sample. The reads are then
aligned to a transcriptome, counted per gene and normalised between samples \cite{stark2019}. Since its
introduction \cite{wang2009}, RNA-seq has become the de facto method for measuring transcriptome-wide
expression, replacing hybridisation microarrays because it is cheaper and easier to implement and, above
all, because it yields comprehensive, quantitative data across the transcriptome \cite{stark2019}. A bulk
library needs microgram amounts of total RNA, which corresponds to hundreds of thousands of mammalian
cells \cite{tang2009}, so each gene is measured as an aggregate over all of them, and differences between
the cells are masked (Section~\ref{sec:bg-heterogeneity}).

scRNA-seq measures each cell separately. An early assay applied mRNA-seq to
single mouse blastomeres \cite{tang2009}. Cells are isolated in wells or droplets,
and their messenger RNA is captured, reverse-transcribed and amplified. Every molecule is labelled with a
cellular barcode that identifies its well or droplet and, in many protocols, with a unique molecular
identifier (UMI) that distinguishes amplified copies of one molecule from separate molecules of the same
gene \cite{luecken2019}. The droplet system of 10x Genomics, on which the SCP542 atlas was profiled
(Appendix~\ref{app:qc}), encapsulates cells with barcoded gel beads and counts the 3$'$ ends of
messenger RNA in tens of thousands of cells per sample \cite{zheng2017}. Droplet systems profile the
most cells, whereas plate-based methods capture fewer cells but detect more genes per cell
\cite{stark2019}. The result is a matrix of counts, barcodes by genes, and it is far sparser than bulk
data. Only 10--20\,\% of a cell's transcripts are reverse-transcribed \cite{stark2019}, so the count
depth differs between identical cells and many entries are zero \cite{luecken2019}. These zeros were long
attributed to a separate technical failure, termed dropout, but in droplet data with UMIs their number
matches what the sampling of molecules predicts, and zeros beyond that reflect biological variation
\cite{svensson2020}. A barcode may also carry two cells (a doublet) or none \cite{luecken2019}.
Normalisation for count depth (Section~\ref{sec:bg-representation}) and quality control
(Section~\ref{sec:data}) address the first and the last of these.

To profile many cancer cell lines simultaneously, Kinker et al.\ pooled lines from diverse lineages,
sequenced each pool together and assigned each cell to its line afterwards, by single-nucleotide polymorphisms
detected in the reads and by mapping expression clusters to the line with the most similar bulk RNA-seq
profile \cite{scp542}. The result is single-cell expression for cell lines that also carry bulk drug-response
labels, which is the pairing this work rests on: the model reads single-cell profiles, while the labels
it is trained on remain bulk measurements (Section~\ref{sec:bg-screens}).

\section{Representing single-cell data}
\label{sec:bg-representation}

A count matrix is rarely analysed as it comes. A standard analysis first normalises the counts, then
reduces their dimension, and only then clusters, visualises or models the cells
\cite{luecken2019,wolf2018scanpy}. The most common normalisation scales each cell's counts by its count
depth, to counts per million (CPM) or a similar target. This assumes that all cells initially held the same
number of messenger RNA molecules and that differences in depth arise only from sampling
\cite{luecken2019}. The normalised values are then typically transformed as $\log(x+1)$, which has three
effects: distances become log fold changes, the dependence of a gene's variance on its mean is mitigated
though not removed, and the skew of the data is reduced towards the normal distribution that many
downstream methods assume \cite{luecken2019}. Variation that normalisation leaves in place, such as batch
effects between cells handled in distinct groups, may need separate correction. Pooling cells from
several samples and assigning them afterwards, for example by genetic variation, can avoid batch effects
between those samples \cite{luecken2019}, and the SCP542 lines were profiled in such pools
(Section~\ref{sec:bg-rnaseq}).

Dimension is reduced in two steps. Feature selection keeps the genes that are informative of the
variability in the data, the highly variable genes (HVGs), typically between 1,000 and 5,000 of them
\cite{luecken2019}. A widely used criterion computes each gene's mean and dispersion (variance over mean),
groups genes into bins by mean expression and standardises the dispersion within each bin, so that a gene
counts as variable relative to genes of similar abundance \cite{satija2015,wolf2018scanpy}. Principal
component analysis then creates new, uncorrelated variables that successively maximise the variance
retained. They are found by solving an eigenvalue problem and are defined by the dataset at hand rather
than fixed in advance \cite{jolliffe2016}. PCA is linear and captures the structure of the data in few
dimensions less well than non-linear methods, but it is the basis of many tools for clustering and
further analysis \cite{luecken2019}. A representation learned in advance on many other datasets is the
alternative that single-cell foundation models offer (Section~\ref{sec:bg-foundation}).

For display, a common workflow reduces the data first to a few dozen principal components and then
non-linearly to two dimensions \cite{chari2023}. Uniform manifold approximation and projection (UMAP) is
such a method. It is competitive with t-distributed stochastic neighbour embedding (t-SNE) for
visualisation and arguably preserves more of the global structure, at lower run time
\cite{mcinnes2018umap}, and it has been regarded as best practice for exploratory visualisation of
single-cell data \cite{luecken2019}. Reducing hundreds or thousands of dimensions to two, however,
inevitably induces substantial distortion, and Chari and Pachter argue that such embeddings are
counter-productive for exploratory biological analysis \cite{chari2023}. They remain useful for showing
structure, but distances in them are not measurements, so a quantitative claim here rests on the full
representation and a two-dimensional map only illustrates it.

\section{Cancer cell-line drug screens and dose--response measures}
\label{sec:bg-screens}

Cancer cell lines derived from tumours recapitulate, in many respects, the tissue type and genomic
context of cancer, and they are a standard tool in molecular biology and drug discovery
\cite{gdsc_yang2013}. The Cancer Cell Line Encyclopedia (CCLE) compiled gene expression, copy number and
sequencing data for 947 human cancer cell lines and, with pharmacological profiles of 24 drugs across 479
of them, identified genetic, lineage and expression-based predictors of drug sensitivity
\cite{barretina2012}. Its later characterisation extended to 1,072 lines and added RNA-seq for 1,019 of
them \cite{ghandi2019}. These bulk profiles are what most response models take as input
(Section~\ref{sec:bg-prediction}).

Several screens measure the response of such lines at scale. GDSC measures viability with fluorescence-based assays after 72 hours of treatment at nine
concentrations, and at its 2013 description held data for 138 drugs across almost 700 lines
\cite{gdsc_yang2013}. CTRPv2 measured 860 lines against an
informer set of 481 compounds --- approved drugs, clinical candidates and probes --- at 16 concentrations
in duplicate, with cellular ATP after 72 hours as the surrogate for viability \cite{ctrpv2}. PRISM instead labels
cell lines with DNA barcodes and screens them in pools, with the relative abundance of each barcode as
the surrogate for viability. Its repurposing screen tested 4,518 drugs against 578 lines at one dose and
retested 1,448 active drugs at eight doses \cite{corsello2020,prism}. In all of them, the readout comes from a well of many cells, so
the response of a cell line is a population average (Section~\ref{sec:bg-heterogeneity}).

The measurements at each concentration, relative to untreated controls, are fitted to a sigmoidal
dose--response curve and summarised by the concentration at which the response is half that of the
control (IC50), the response at the highest concentration (Emax), or the area under the curve (AUC)
\cite{hafner}. Where a curve does not fall to half of the control within the tested concentrations, the IC50 has to
be extrapolated, as the Cancer Genome Project (CGP), the precursor of GDSC, did with a sigmoid model, or replaced by a placeholder such
as the highest concentration, as CCLE did. The AUC requires no extrapolation and can always be estimated
\cite{haibekains2013} (Section~\ref{sec:target}). All three depend on how
many divisions the cells undergo during the assay, which creates artefactual associations with genotype,
and growth-rate-normalised measures have been proposed to correct for it \cite{hafner}. Fitting itself
can mislead: a curve fitted to noise yields potency and effect estimates like any other, and treating
those with the same confidence as true responses can lead to wrong hypotheses and to problems in
training machine-learning models. CurveCurator therefore fits a log-logistic model and computes, for
each curve, a $p$-value and false discovery rate against a curve without response
\cite{bayer2023curvecurator}.

The screens do not agree well with one another. For the 471 lines and 15 drugs that both CGP and CCLE
assayed, gene expression was well correlated between the
two studies, whereas drug response was highly discordant: only one drug had a Spearman correlation of
measured response above 0.6 \cite{haibekains2013}. DrEval addresses part of this by refitting
every screen it includes from replicate-level data with CurveCurator, so that response measures are
estimated consistently across datasets \cite{dreval}. The labels used here are its refit of CTRPv2
(Section~\ref{sec:target}).

\section{Computational drug-response prediction and its evaluation}
\label{sec:bg-prediction}

Drug-response prediction models learn the response of a cancer sample to a drug from features of the
sample and, in many models, of the drug \cite{partin2023}. Drug-specific models learn from the cancer
features alone, one model per drug. Multi-task models take the same input and produce one output per
drug, which lets them learn from more response data and exploit what drugs have in common
\cite{partin2023}. CCLE already built such predictors with naive Bayes classifiers and elastic-net
regression \cite{barretina2012}. The NCI-DREAM challenge then compared 44 algorithms on breast cancer
lines: the best modelled non-linear relationships and used pathway information, gene expression
microarrays gave the best predictive power of any single profiling data set, and combining data sets
improved performance further \cite{costello2014}. Gene expression has since been used, alone or with
mutation, copy-number and other data, in about 90\,\% of the 61 deep-learning models reviewed by Partin
et al., whose most common targets are the IC50 and the AUC \cite{partin2023}. The task is difficult
partly because of the data: as models grow in complexity, the number of observations they need grows
with them, and deep networks need many samples to avoid overfitting \cite{adam2020}.

How such models are evaluated has proved as consequential as how they are built. DrEval identifies
obstacles that make reported performance overly optimistic \cite{dreval}. First, the split must match
the intended use: a model meant for unseen patients must be tested on cell lines not used for training
or tuning (leave-cell-line-out, LCO), whereas splitting random pairs of line and drug (leave-pairs-out, LPO) is warranted
only for imputing missing values. Second, the measurements are pseudo-replicates: CTRPv2's roughly
400{,}000 response curves in DrEval's release stem from \NCtrpLines{} cell lines and \NDrugs{} compounds, so the effective
sample size is far smaller than the number of curves. Third, most of the variance in response lies
between drugs, because drugs act at very different concentrations, so a model that memorises each drug's
mean explains most of it. This Simpson's paradox inflates global metrics. In one model, a global
Pearson correlation of 0.91 fell to 0.56 when computed per drug. Evaluated on these terms, deep-learning
models barely outperform a naive model that predicts only the mean drug and cell-line effects, about half
of the tested models do not significantly outperform it, and under LCO no model surpasses a tuned random
forest \cite{dreval}. This work adopts those terms: LCO splits, scores computed within each compound, and
DrEval's naive baselines (Section~\ref{sec:evaluation}).

\section{Drug response at single-cell resolution}
\label{sec:bg-singlecell}

Single-cell chemical screens exist, but they measure a different quantity. sci-Plex was introduced as an
alternative to coarse readouts such as cell survival. It quantifies the global transcriptional response
to thousands of perturbations at single-cell resolution, and its first screen profiled about 650{,}000
cells of three cancer cell lines exposed to 188 compounds \cite{srivatsan2020}. Tahoe-100M measures how
1,100 small molecules affect 100 million cells across 50 cancer cell lines \cite{tahoe100m}. Their
readout is the transcriptome after treatment, which is not a viability label of the kind the screens of
Section~\ref{sec:bg-screens} provide. Methods that predict response per cell therefore combine
single-cell expression with bulk labels, and they do so in different ways.

Beyondcell derives drug sensitivity signatures from cell-line screens, by differential expression
between lines with high and low AUC, scores each cell for its enrichment in these signatures and groups
cells into therapeutic clusters with distinct predicted responses \cite{beyondcell2021}. scDEAL trains on
bulk expression and response of cell lines, learns a feature space shared between bulk and single-cell
data, and transfers the model to individual cells \cite{scdeal}. Its authors name the limited amount of
benchmark data as the main obstacle. They evaluated it on six scRNA-seq datasets treated with five drugs,
whose per-cell labels are binary and, in most of the original studies, were assigned to whole groups of
cells by their treatment condition \cite{scdeal}.

PERCEPTION is the closest to this work, because it uses the same atlas \cite{perception,scp542}. For
each drug it fits an elastic-net model on bulk expression and PRISM response of cell lines and then
applies it to single-cell expression, tuning the number of genes on the lines that have both. A line's
prediction is the mean over its cells, scored with leave-one-line-out cross-validation. Predictive
models were obtained for 44 of 133 drugs, and a model built on single-cell expression alone, without
pretraining on bulk data, reached a Pearson correlation of 0.22 against 0.39 on held-out lines
\cite{perception}. The present work trains on single cells directly, with the bulk label broadcast to
every cell of a line and no bulk pretraining. In PERCEPTION's linear models, leaving out that
pretraining gave the weaker result. This work asks
whether a pretrained single-cell representation can take the place of that pretraining
(Section~\ref{sec:bg-foundation}).

\section{Single-cell foundation models}
\label{sec:bg-foundation}

A transformer processes a set of tokens with attention mechanisms alone, without recurrence or
convolution \cite{vaswani2017}. Each token is a vector, and attention replaces it with a weighted sum over
the values of all tokens, where the weight is computed from the similarity of its query to their keys. In
scaled dot-product attention, with the queries, keys and values of all tokens stacked into matrices $Q$,
$K$ and $V$ and keys of dimension $d_k$,
\begin{equation}
  \mathrm{Attention}(Q, K, V) = \mathrm{softmax}\!\left(\frac{QK^\top}{\sqrt{d_k}}\right) V
  \label{eq:attention}
\end{equation}
\cite{vaswani2017}. Several such attention functions run in parallel as multi-head attention, and
transformer layers are stacked. The representation of each token therefore depends on every other token
in the input.

Single-cell foundation models apply this architecture to transcriptomes. They follow the transfer
learning that changed natural language processing and computer vision: a model is pretrained on a large,
general dataset and then fine-tuned towards downstream tasks that have limited task-specific data
\cite{theodoris2023}. Geneformer was pretrained on 29.9 million human single-cell transcriptomes from a
broad range of tissues. It represents each cell by its genes ranked by expression and learns with a
self-supervised masked objective \cite{theodoris2023}. scGPT was pretrained on 33 million human cells
under normal, non-disease conditions from the CELLxGENE collection \cite{scgpt}. It treats genes as tokens,
converts each cell's expression into relative values by binning, and produces cell and gene embeddings
together. Fine-tuned, it reached strong performance in cell type annotation, multi-batch integration,
perturbation prediction and gene network inference \cite{scgpt}.

Such models can also be used without fine-tuning, with the embedding taken as it is (zero-shot). In that
setting their advantage has not held. Separating known cell types across five datasets, scGPT and
Geneformer performed worse than a selection of highly variable genes and the established methods Harmony
and scVI, and the highly variable genes outperformed both foundation models on all metrics, although the
datasets resembled those used in pretraining \cite{kedzierska2025}. The authors of scGPT note that its
current pretraining does not inherently mitigate batch effects, so its zero-shot performance may be
limited on data with substantial technical variation \cite{scgpt}.

For this work, scGPT offers a representation learned from far more cells than any drug screen labels,
and one that was not fitted to the atlas it is applied to. Against that, its pretraining corpus contains
no disease cells, so cancer cell lines were not among them, and its zero-shot embeddings have been
weaker than highly variable genes on other tasks. Whether its frozen embedding is the better input for
predicting drug response is therefore a question to test, not a premise (Section~\ref{sec:representations}).

\section{Multiple-instance learning}
\label{sec:bg-mil}

In multiple-instance learning (MIL), a label is attached not to a single instance but to a bag of
instances, which have no order among them and whose own labels exist but are unknown during training
\cite{ilse2018}. The problem was first posed for drug activity. A molecule can adopt several shapes, so a
bag holds the conformations of one molecule, and the molecule has the desired effect if and only if at
least one of them binds the target \cite{dietterich1997,ilse2018}. In the standard formulation the bag
label is therefore binary and equals the maximum of the instance labels, $Y = \max_k y_k$, so a model
must be invariant to the order of the instances \cite{ilse2018}.

A MIL model consists of a transformation of each instance, a permutation-invariant pooling over the
bag, and a final transformation to the bag label. In the instance-level approach every instance receives
its own score and the scores are pooled. In the embedding-level approach the instances are pooled into
one bag representation first, which is then classified. The embedding-level approach has been argued to
classify bags better, because an instance-level classifier is trained without instance labels and can
add error \cite{ilse2018}. Common pooling operators are the maximum and the mean. Ilse et al.\ replace
them with a learned weighted average of the instance embeddings $h_k$,
\begin{equation}
  z = \sum_{k=1}^{K} a_k h_k, \qquad
  a_k = \frac{\exp\{w^\top \tanh(V h_k^\top)\}}{\sum_{j=1}^{K} \exp\{w^\top \tanh(V h_j^\top)\}},
  \label{eq:attention-mil}
\end{equation}
whose weights show how much each instance contributes to the bag label \cite{ilse2018}.

A cell line fits this frame loosely. Its cells are the instances, the line is the bag, and the drug
response is measured only for the bag, while each cell's own response remains unknown
(Section~\ref{sec:bg-screens}). Unlike the standard case, however, the label is continuous rather than
binary. How this work adapts MIL to such a label is described in Section~\ref{sec:mil}.

\section{Learning and evaluation methods}
\label{sec:bg-methods}

The per-cell model is a multilayer perceptron, a stack of linear layers with non-linear activations, here
the Gaussian error linear unit (GELU) \cite{hendrycks2016gelu}. Layer normalisation rescales the
activations within each layer \cite{ba2016layernorm}, and dropout sets a random share of units to zero
during training to reduce overfitting \cite{srivastava2014}. Such networks are trained by stochastic
gradient descent, commonly with Adam, which scales each parameter's step by running estimates of the
gradient's moments \cite{kingma2015adam}, or AdamW, which applies weight decay separately from that
scaling \cite{loshchilov2019adamw}. Clipping the gradient norm prevents single steps from becoming too
large \cite{pascanu2013}, and early stopping ends training when the error on held-out data stops
improving \cite{prechelt1998}. Initialising an output's bias at the target's mean saves the first epochs
from fitting that level \cite{lin2017focal}.

When continuous targets are unevenly distributed, a regression model fits the crowded middle and neglects
rare values. Weighting each sample by the inverse of the target's density, estimated for example by a
kernel density estimate \cite{scott1992}, counteracts this \cite{dir,denseweight}.

Simple models serve as controls. Ridge regression adds a squared penalty on the coefficients to a linear
least-squares fit \cite{hoerl1970}, and its penalty can be chosen by leave-one-out error computed in
closed form \cite{hastie2009,pedregosa2011}. A random forest averages many decision trees, each grown on a
resample of the data \cite{breiman2001}.

Rankings are compared with the Spearman correlation, the Pearson correlation of ranks
\cite{spearman1904}. How well a score separates two groups is measured by the area under the receiver
operating characteristic curve (AUROC), which equals the Mann--Whitney statistic divided by the product of
the group sizes \cite{mannwhitney1947,hanley1982}. The uncertainty of an estimate can be obtained by the
bootstrap, recomputing it on resamples of the independent units \cite{efron1993}, and the calibration of
a model's predictions by the slope of a regression of truth on prediction \cite{vancalster2016}.

\section{Objective of this work}
\label{sec:bg-objective}

Drug response is measured on populations of cells, whereas resistance can arise from subpopulations
that a population measurement averages away (Section~\ref{sec:bg-heterogeneity}). A model that predicts
response for individual cells could report the distribution of sensitivities within a sample and, in
principle, point to resistant cells before treatment selects for them. The long-term goal of this
project is a pan-cancer single-cell foundation model of pharmacological response, for both efficacy and
toxicity, that can later be fine-tuned to specific cancer types or to clinical data with binary
outcomes. Such a model has to be trained on labels that exist at scale, and those labels are bulk
measurements on cell lines (Section~\ref{sec:bg-screens}).

The objective of this work is the first step towards that goal. It is to establish, on one atlas and one
screen, whether and how well drug efficacy can be learned from single cells when every label is a bulk
measurement, and what such a model does and does not capture. Three questions follow from it.
\begin{enumerate}
  \item Can drug response be predicted from single cells with bulk labels at all?
  \item Are scGPT embeddings a viable representation for this task, measured against principal component
        analysis at matched width?
  \item Does a model trained on single cells learn cellular heterogeneity implicitly?
\end{enumerate}
The first question decides whether the setting supports a model at all. The second asks whether a
pretrained single-cell representation is a better starting point than one fitted to the data at hand,
which is the premise of the foundation-model approach (Section~\ref{sec:bg-foundation}). The third is
the clinically relevant one. A model that reproduces only the population mean cannot see a surviving
subpopulation, however accurate that mean is, so the value of working at single-cell resolution rests
on its answer.
